% Generated from lessons/0002-cluster-closure-boundary.html; edit the HTML source, then rerun the converter.
\section{聚点、闭包、边界}
\lessonmark{聚点、闭包、边界}

这节课解决一个常见卡点：点在不在集合里，和它是不是聚点、闭包点、边界点，不是一回事。考试题常常就考这层区别。

\subsection*{先分清三句话}

\begin{conceptbox}
\textbf{闭包点}
每个小球都碰到 \(S\)。允许只碰到点自己。

\[
          B(x,r)\cap S\ne\varnothing
        \]
\end{conceptbox}

\begin{conceptbox}
\textbf{聚点}
每个小球都碰到 \(S\) 中别的点。不能只靠 \(x\) 自己。

\[
          B(x,r)\cap(S\setminus\{x\})\ne\varnothing
        \]
\end{conceptbox}

\begin{conceptbox}
\textbf{边界点}
每个小球既碰到 \(S\)，也碰到 \(S\) 的外面。

\[
          B(x,r)\cap S\ne\varnothing,\qquad B(x,r)\cap S^c\ne\varnothing
        \]
\end{conceptbox}

教材对应：《数学分析》第三版下册第十一章 §1“开集与闭集”“聚集”。教材给出聚点的序列判别：\(x\) 是 \(S\) 的聚点，当且仅当存在 \(S\) 中互异点列 \(x_k\to x\)。

\subsection*{最容易混的区别}

\begin{center}
\small
\begin{tabularx}{\linewidth}{|>{\raggedright\arraybackslash}p{0.18\linewidth}|>{\raggedright\arraybackslash}p{0.42\linewidth}|>{\raggedright\arraybackslash}X|}
\hline
\textbf{问题} & \textbf{答案} & \textbf{例子} \\ \hline
聚点一定属于集合吗？ & 不一定。 & \(0\) 是 \(\{1/n:n=1,2,\ldots\}\) 的聚点，但 \(0\) 不在这个集合中。 \\ \hline
集合中的点一定是聚点吗？ & 不一定。 & \(1\) 是 \(\{1/n\}\) 中的点，但它是孤立点，不是聚点。 \\ \hline
边界点一定属于集合吗？ & 不一定。 & \((0,1)\) 的边界点 \(0\) 和 \(1\) 都不属于 \((0,1)\)。 \\ \hline
闭包是什么？ & 所有贴得上 \(S\) 的点。 & \(\overline{(0,1)}=[0,1]\)；\(\overline{\{1/n\}}=\{1/n:n\ge1\}\cup\{0\}\)。 \\ \hline
\end{tabularx}
\end{center}

期末判断时，先问“小球碰到谁”：只碰 \(S\) 是闭包点；碰 \(S\) 中别的点是聚点；同时碰 \(S\) 和补集是边界点。

\subsection*{例题：\(S=\{1/n:n=1,2,\ldots\}\)}

这个例子很适合拆概念。

\[
      S=\{1,1/2,1/3,\ldots\}\subset R
    \]

\(0\) 是聚点，因为任意小区间 \((-r,r)\) 中，总能找到足够大的 \(n\)，使 \(1/n<r\)。

但每个 \(1/n\) 都是孤立点。比如点 \(1/3\)，左右离最近的其他点还有正距离，所以能取一个很小的邻域，里面只有 \(1/3\) 这一个 \(S\) 中的点。

因此：

\[
      S'=\{0\},\qquad \overline S=S\cup\{0\},\qquad \partial S=S\cup\{0\}
    \]

\subsection*{马上练}

判断给出的点属于哪一类。选完立刻看反馈。

\begin{quickcheck}
\(S=\{1/n:n=1,2,\ldots\}\)。点 \(x=0\) 是 \(S\) 的什么点？

\tcblower
\textbf{答案：}聚集点。任意小区间里都有 \(1/n\) 这样的点靠近 \(0\)，而且这些点都不等于 \(0\)。\(0\) 不在 \(S\) 中不影响它成为聚点。
\end{quickcheck}

\begin{quickcheck}
\(S=\{1/n:n=1,2,\ldots\}\)。点 \(x=1\) 是 \(S\) 的什么点？

\tcblower
\textbf{答案：}孤立点。点 \(1\) 属于 \(S\)，但在足够小的邻域里没有 \(S\) 中其他点，所以它不是聚点，而是孤立点。
\end{quickcheck}

\begin{quickcheck}
\(S=(0,1)\)。点 \(x=1\) 是 \(S\) 的什么点？

\tcblower
\textbf{答案：}边界点。以 \(1\) 为中心的任意小区间都同时碰到 \((0,1)\) 和它的补集；\(1\) 不属于集合也仍然可以是边界点。
\end{quickcheck}

\begin{quickcheck}
\(S=(0,1)\)。点 \(x=1/2\) 是 \(S\) 的什么点？

\tcblower
\textbf{答案：}内部点。\(x=1/2\) 离左右边界都有正距离，可以放一个小区间完全留在 \((0,1)\) 中。内点当然也是闭包点。
\end{quickcheck}

\subsection*{本节收束}

今天的核心动作是：不要先问“这个点在不在集合里”，先问“任意小球会碰到哪些点”。碰到 \(S\)，是闭包语言；碰到 \(S\) 中别的点，是聚点语言；两边都碰到，是边界语言。

整节复习页：第十一章 §1：Euclid 空间上的基本定理。上一课：用球邻域判断开闭。速查表：欧氏空间期末速查表。教材路线：Euclid 空间上的极限和连续。

推荐阅读：本地教材《数学分析》第三版下册第十一章 §1 的“聚集”部分；辅助参考 \href{https://www.math.ucdavis.edu/~hunter/intro_analysis_pdf/ch13.pdf}{UC Davis Chapter 13}。

如果你愿意，下一步把教材第 106 页习题 4 或 5 发给我，我们可以按期末标准一起写一遍。
